\documentclass{report}
\usepackage{amsmath,amssymb}
\setlength{\parindent}{0mm}
\setlength{\parskip}{1em}
\begin{document}
\begin{center}
\rule{6in}{1pt} \
{\large Lei Tang \\
{\bf Parallel FOSLS-AMG for MHD Equations}}

1475 Folsom St Apt 384 \\ Boulder \\ CO 80302
\\
{\tt tangl@colorado.edu}\\
Marian Brezina\\
Jose Garcia\\
Tom Manteuffel\\
Steve McCormick\\
John Ruge\end{center}

We describe the use of an efficiency-based adaptive mesh refinement (AMR)
scheme on a 2D reduced model of the incompressible, resistive
magnetohydrodynamic (MHD) equations. The fully implicit 2-step BDF scheme
is considered for time discretization. At each time step, a First-order
System Least Squares (FOSLS) finite element formulation and algebraic
multigrid (AMG) methods are used in the context of nested iteration. The
AMR scheme chooses elements to refine on the basis of minimizing the
`accuracy-per-computational-cost' efficiency (ACE) measure that takes
into account both error reduction and computational cost. Accommodations
of the ACE approach to massively distributed memory architectures involve
binning strategy to reduce communication cost. Load balancing begins at
very coarse levels and is guided by a space filling curve. To evaluate
solutions at previous time steps on the current mesh, we employ the
quad-tree(forest) structure associated with the mesh at each time step to
reduce communication cost. Numerical results are presented for the
simulation of a large aspect-ratio tokamak instability. We show that
these methods are able to resolve the physics using a lot less
computational cost than uniformly refined mesh to approximate the
solution within the same error tolerance. Weak and strong scalability are
demonstrated up to thousands of processors on an IBM Blue Gene/L
supercomputer.


\end{document}
